\section{Agent 的定义与未来}

\subsection{Agent 的两个核心特征}

\begin{importantbox}{Agent = 多轮 + 工具}
杨植麟对 Agent 的定义简洁而本质：
\begin{enumerate}
    \item \textbf{多轮}（Multi-turn）：可以进行多次交互，这是 Test-Time Scaling 的一种方式，使模型能完成复杂任务
    \item \textbf{工具}（Tool Use）：连接大脑与外部世界的方式——搜索引擎连接互联网，代码能力连接数字世界的一切自动化
\end{enumerate}
\end{importantbox}

\subsection{通用性是 Agent 最缺的能力}

\begin{importantbox}{从通用 Agent 到垂直应用}
如果 Agent 有足够好的泛化能力，就不需要大量的垂直 Agent：
\begin{itemize}
    \item 通用 Agent 只需接入不同的工具（定制数据库、定制 API、定制文档接口等），就能完成垂直领域任务
    \item 但当前 Agent 最缺的正是这种泛化能力——能否泛化到训练中没有见过的工具和环境
    \item Benchmark 过拟合是一个风险：刷几个 Benchmark 的分数，不代表真实场景的泛化能力
\end{itemize}
\end{importantbox}

\begin{knowledgebox}{Agent 的目的不是模拟人}
杨植麟指出一个重要的区分：Agent 的目的是通用（General Purpose），而不是模拟人。人类也是通用的（Universal Constructor），Agent 的行为方式与人类相似只是一个巧合结果——就像飞机的设计目的是交通工具，不是为了像鸟一样飞。
\end{knowledgebox}

\subsection{本章小结}

Agent 的两个核心特征是多轮交互和工具使用。当前最大挑战是泛化能力——能否在没见过的工具和环境上表现良好。Agent 的目标是通用智能，不是模拟人类行为。


